% 第1章 金融取引と金融市場(資料1)
\chapter{金融取引と金融市場}

本章では，本講義全体の出発点として，金融 (finance) とは何か，経済主体はどのように
資金を調達するのか，そして証券市場 (securities market) がどのような役割を果たして
いるのかを概観する。数理的な議論は次章から始まるが，そこで扱う債券・株式などの
金融商品や市場の仕組みを，まずことばとして正確に理解しておくことが本章の目的である。

\section{金融}

\subsection{金融とは何か}

\begin{definition}[金融]
\textbf{金融}とは，余剰資金のある経済主体 A から，資金の不足している経済主体 B へ
お金を融通することをいう。
\end{definition}

このとき，資金を出す側の A を\textbf{貸し手}，資金を受け取る側の B を\textbf{借り手}と
呼ぶ。A の側から見ればこの行為は\textbf{投資}であり，B の側から見れば\textbf{調達}である。
すなわち，同じ資金の流れを，貸し手の視点で見るか借り手の視点で見るかによって
呼び方が変わる。

用語をいくつか整理しておく。貸し借りされるお金の本体を\textbf{元本}という。
元本の使用の対価として，借り手が支払うお金を\textbf{利子}，貸し手が受け取るお金を
\textbf{利息}と呼び分けることがある(講義では厳密には区別しないが，このような
使い分けがあることは知っておくとよい)。

\subsection{直接金融と間接金融}

資金の融通のされ方には，\textbf{直接金融}と\textbf{間接金融}の2つの形態がある。
両者の違いは，仲介者がどのような立場で取引に関与するかにある。

\begin{itemize}
\item \textbf{直接金融}：経済主体 A が，証券会社等の仲介者を通して，経済主体 B に
\emph{直接}資金を貸す形態。元本は A から B へ渡り，利子は B から A へ支払われる。
仲介者(証券会社等)は取引を仲介するだけであり，その対価として B から
\textbf{手数料}を受け取る。
\item \textbf{間接金融}：経済主体 A がまず仲介者(銀行等)に資金を預け，仲介者が
自らの名義で経済主体 B に貸し出す形態。元本は A $\to$ 銀行 $\to$ B と流れ，
利子は B が銀行に利率 $g$ で支払い，銀行は A に利率 $f$ で支払う。ここで
$f < g$ であり，その差(利鞘)が銀行の収益となる。
\end{itemize}

\begin{figure}[htbp]
\centering
\begin{tikzpicture}[
  box/.style={draw, rectangle, minimum width=2.4cm, minimum height=1.0cm, align=center},
  arr/.style={-{Stealth}, thick},
  node distance=3.2cm]
% 直接金融
\node[box] (A1) {経済主体\\A};
\node[box, right=of A1] (M1) {仲介者\\(証券会社等)};
\node[box, right=of M1] (B1) {経済主体\\B};
\draw[arr, red!70!black] ([yshift=2mm]A1.east) -- ([yshift=2mm]B1.west)
  node[pos=0.75, above]{元本};
\draw[arr, red!70!black] ([yshift=-2mm]B1.west) -- ([yshift=-2mm]A1.east)
  node[pos=0.75, below=1mm]{利子(利息)};
\draw[arr, blue!70!black] ([yshift=-4mm]B1.west) -- ([yshift=-4mm]M1.east)
  node[pos=0.5, below]{手数料};
\node[above=1mm of M1.north] {\textbf{直接金融}};
\end{tikzpicture}

\vspace{6mm}

\begin{tikzpicture}[
  box/.style={draw, rectangle, minimum width=2.4cm, minimum height=1.0cm, align=center},
  arr/.style={-{Stealth}, thick},
  node distance=3.2cm]
% 間接金融
\node[box] (A2) {経済主体\\A};
\node[box, right=of A2] (M2) {仲介者\\(銀行等)};
\node[box, right=of M2] (B2) {経済主体\\B};
\draw[arr, red!70!black] ([yshift=2mm]A2.east) -- ([yshift=2mm]M2.west)
  node[pos=0.5, above]{元本};
\draw[arr, red!70!black] ([yshift=-2mm]M2.west) -- ([yshift=-2mm]A2.east)
  node[pos=0.5, below]{利子(利息) $f$};
\draw[arr, red!70!black] ([yshift=2mm]M2.east) -- ([yshift=2mm]B2.west)
  node[pos=0.5, above]{元本};
\draw[arr, red!70!black] ([yshift=-2mm]B2.west) -- ([yshift=-2mm]M2.east)
  node[pos=0.5, below]{利子(利息) $g$};
\node[above=1mm of M2.north] {\textbf{間接金融}};
\node[below=5mm of M2.south, blue!70!black] {$f < g$};
\end{tikzpicture}
\caption{直接金融と間接金融における資金の流れ。直接と間接では，元本と利子の
「起点と終点」が異なることに注意する。}
\label{fig:direct-indirect}
\end{figure}

図\ref{fig:direct-indirect}のとおり，直接金融では元本・利子の流れの起点と終点は
A と B であるのに対し，間接金融ではいったん銀行を経由するため，A と B の間に
直接の貸借関係はない。この違いは，借り手が返済不能になったときに誰が損失を
被るかという点で決定的に重要である。

\begin{keypoint}[借り手 B が資金繰りに行き詰まったら?]
\begin{itemize}
\item \textbf{直接金融の場合}：貸しているのは経済主体 A 自身なので，損失を被るのは
A である。仲介者である証券会社は原則として損失を負わない。
\item \textbf{間接金融の場合}：B に貸しているのは銀行なので，まず損失を被るのは
銀行である。A は銀行に預金しているだけなので，銀行が健全である限り A の資金は
守られる。
\end{itemize}
つまり，直接金融では投資家がリスクを直接負担し，間接金融では銀行がリスクを
吸収する構造になっている。間接金融で $f<g$ という利鞘が存在するのは，銀行が
このリスク負担と仲介機能の対価を得ているためと解釈できる。
\end{keypoint}

\subsection{投資家，金融機関，事業会社}

金融取引に登場する経済主体は，次のように分類できる。

\begin{description}
\item[投資家] お金を融通する側の経済主体。
  \begin{itemize}
  \item \textbf{個人投資家}：個人としてお金を融通している経済主体。
  \item \textbf{機関投資家}：企業としてお金を融通している経済主体。顧客から集めた
  資金を運用・管理する法人投資家であり，保険会社，投資信託会社，信託銀行，
  年金基金，投資顧問会社，ヘッジファンドなどが該当する。
  \end{itemize}
\item[金融機関] 金融の仲介を業とする主体。
  \begin{itemize}
  \item \textbf{証券会社}：直接金融を営む金融機関。
  \item 間接金融を営む金融機関は多数ある。例として，普通銀行，郵便貯金銀行，
  信託銀行，信用金庫，信用協同組合などが挙げられる。
  \end{itemize}
\item[(一般)事業会社] 金融業以外の一般の企業。通常はお金の借り手側に立つ。
\end{description}

\section{資金調達}

事業会社が資金を調達する手段は，直接金融によるもの(有価証券の発行)と
間接金融によるもの(銀行借入)に大別される。

\subsection{直接金融における資金調達：有価証券}

直接金融では，証券会社を通して\textbf{有価証券}を取引することで資金のやり取りが
行われる。

\begin{definition}[有価証券]
\textbf{有価証券}とは，財産権を表章したものであり，それに記載された権利の行使や
移転(譲渡)がその証券によって行われるものをいう。
\end{definition}

有価証券は\textbf{債券}と\textbf{株券}に大別される。両者は資金調達者の貸借対照表上の
位置づけがまったく異なる。

\begin{itemize}
\item 債券で調達される資金は\textbf{負債}に該当する。負債とは返済義務のある契約，
すなわち借りているお金のことである。
\item 株券で調達される資金は\textbf{株式}と呼ばれ，自己資本に相当する。
\end{itemize}

以下で見るように，債券と株式は性質が大きく異なる。

\subsubsection{債券}

\begin{definition}[債券]
\textbf{債券} (bond) とは，資金調達を行う発行者(\textbf{発行体})が，お金を借りた
証拠として，利息の支払いや元本の返済を約束して発行する証書である。
\end{definition}

債券の保有者が受け取る将来のキャッシュフローは，\emph{ある程度は確定的}である。
すなわち，決まった日に決まった金額を受け取る(ただし発行体の破綻などの例外は
ある)。この「キャッシュフローの確定性」が，次章以降で債券を数理的に扱う際の
出発点となる。

債券に関する基本用語を整理する。
\begin{itemize}
\item \textbf{債権者}：お金を貸している側，すなわち投資家。
\item \textbf{債務者}：お金を借りている側，すなわち発行体。
\item \textbf{額面}：償還時に返済される元本の金額。
\item \textbf{利払い(クーポン)}：定期的に支払われる利息。
\item \textbf{償還日}：元本が返済される日(満期)。
\item \textbf{利払日}：クーポンが支払われる日。
\item \textbf{クーポンレート}：額面に対するクーポンの割合(年率で表示)。
\end{itemize}
日本では通常，クーポンは年2回払い，固定利率で，額面は1億円単位である。

\subsubsection{株式}

株式は債券とは対照的な性質をもつ。

\begin{itemize}
\item 投資家への\textbf{返済義務なし}，\textbf{満期なし}，\textbf{確定的な利払なし}。
\item 株式は自己資本に相当するので，\textbf{株主は当該事業会社の所有者}である。
\item 株主が受け取るお金はクーポンではなく\textbf{配当}である。配当金は業績次第で
決まり，業績によっては無配当になることもある。
\end{itemize}

株式投資家の将来キャッシュフローについては次の点が重要である。
\begin{itemize}
\item 投資の狙いは，\textbf{配当}と，\textbf{キャピタルゲイン}($\approx$ 売買益)を
得ることである。
\item 配当は定額ではなく，受け取れないこともあるため，不確実性が高い。
\item その結果，債券に比べると，株式の価値の変動性は相対的に高い。
\end{itemize}

また，株主は「会社の所有者」であるから，株主総会への参加権と議決権をもつ。
会社経営は経営陣が行うが，重要案件は株主総会を通す必要がある。議決権は株式の
保有枚数に応じて与えられるため，大株主ほど強い発言力をもつ。

\subsection{間接金融における資金調達：銀行借入}

\textbf{銀行借入}は間接金融による資金調達手段である。銀行借入による資金は，
事業会社にとって\textbf{負債}に該当する。この点で債券に似ている(資金調達者である
事業会社にとって，どちらも負債である)が，間接金融であることに注意が必要である。
本講義では，間接金融の貸し手を総じて\textbf{銀行}と呼ぶことにする。

\begin{itemize}
\item \textbf{融資}：銀行が事業会社に資金を貸すこと。
\item 用語：債権者(銀行)，債務者(事業会社)，元本，利払い，満期日，利払日，利率。
\end{itemize}

\subsubsection{預金：銀行の資金調達}

銀行が貸し出す資金は，投資家(預金者)から集めたものである。

\begin{itemize}
\item \textbf{預金}：投資家が銀行に貸し出した資金。特に，ゆうちょ銀行への預金は
\textbf{貯金}という。
\item \textbf{定期性預金}：満期がある預金。満期前に口座から引き出すと，ペナルティー
(預入期間の利息の減額)が科せられる。
\item \textbf{普通預金}：満期がなく，いつでも口座から引き出せる預金
(\textbf{要求払預金}とも呼ばれる)。
\end{itemize}

利率の大小関係について，次の点を押さえておく。事業会社への融資の利率に比べると，
預金の利率はかなり低い。さらに，普通預金の利率は定期性預金の利率よりも低い。
いつでも引き出せるという流動性の高さの代償として，利率が低くなっていると
解釈できる。なお，要求払預金には利息のつくものとつかないものがあり，事業会社の
決済用の\textbf{当座預金}には利息がつかない。

\section{証券市場}

\subsection{探索コストの低減}

金融取引は，余剰資金を貸したい人と，不足資金を借りたい人が\emph{出会わなければ}
成立しない。もし，貸したい人が借りたい人を探索するコスト，あるいは借りたい人が
貸したい人を探索するコストが非常に高かったら，金融取引はほとんど成立せず，
経済活動は大きく阻害されるだろう。したがって，\textbf{探索コストが低いことは，
経済全体にとって重要}である。

\begin{keypoint}[市場と銀行の役割]
直接金融において探索コストを低減し，取引を活発にする役割を果たしているのが
\textbf{証券市場}である。間接金融において同様の役割を果たしているのが
\textbf{銀行}である。
\end{keypoint}

\subsection{証券市場の形態}

\textbf{証券市場}とは，証券を取引する「場所」と考える(ただし実空間上の場所とは
限らない)。取引の形態には，取引所を通す取引と通さない取引がある。

\begin{description}
\item[取引所取引] 日本の(証券)取引所は，東京，大阪，名古屋，札幌，福岡にある。
取引所では定型化された商品(株式，債券など)が取引される。すべての商品を限られた
場所で取り扱うことはできないので，定型化された商品を，大きな取引量で，集中的に
捌くのが取引所の役割である。
\item[店頭取引(店頭市場)] 取引所を通さない取引(市場)。相対取引を軸に，
定型化されていないさまざまな商品，オーダーメイド商品などが取引されている。
\end{description}

市場は取引対象によって，債券市場，株式市場，オプション市場，スワップ市場，
などと呼び分けられる。

\subsection{債券市場}

証券市場は，新規に証券が発行される\textbf{発行市場} (primary market) と，
既に発行された証券が投資家間で売買される\textbf{流通市場} (secondary market) に
分けられる。

債券は，発行体，償還日，クーポンレートなど，何かしらの属性が異なれば別の銘柄と
なるため，\textbf{非常に多くの銘柄}が存在する。

\begin{description}
\item[発行市場] 債券を発行する市場。発行額の規模は，長期国債で1銘柄あたり数兆円，
国債全体では年間約200兆円に達する。通常，額面と(発行)価格は異なり，価格は
額面100円に対する価格で表示される。額面と価格が等しい場合を\textbf{パー} (par) と
いい，パーで発行することを\textbf{パー発行}という。額面より価格が高い場合を
\textbf{オーバーパー}，額面より価格が低い場合を\textbf{アンダーパー}という。
\item[流通市場] 既発債を取引する市場。債券の多くは店頭で取引されるため，実際の
市場価格はつかみにくい。そこで，マーケットメイカーによる売り気配・買い気配の
提示が行われている。価格情報としては，日本証券業協会の売買参考統計値や，
野村證券の JS-Price が有名である。
\end{description}

講義資料では，日本証券業協会のデータに基づく債券の年間発行額の推移が示された。
1998年から2024年頃までの間，債券の年間発行額の合計はおおむね150〜250兆円台の
規模で推移し，その大部分を国債が占める。普通社債の発行額はそれよりはるかに
小さく，年間10〜20兆円程度の規模である。

\subsection{株式市場}

債券と違い，株式では既存の株券も追加発行の株券も同じ扱いであり，ある会社の
株式の価格(\textbf{株価})は1つである。株価は会社の価値を反映するものと
考えられる。

\begin{definition}[株式時価総額]
\[
\text{株式時価総額} = \text{株価} \times \text{発行済株式数}.
\]
\end{definition}

講義資料に示された日本の株式時価総額の実績値を挙げると，東証一部で709兆円・
東証全体で730兆円(2022年3月末)，プライム市場で684兆円・東証全体で712兆円
(2022年4月末)，プライムで970兆円・東証全体で1,008兆円(2024年3月末)，
プライムで911兆円・東証全体で947兆円(2025年3月末)，プライムで1,327兆円・
東証全体で1,372兆円(2026年2月末)となっている。長期の推移を見ると，
1980年代末のバブル期に約600兆円まで拡大した後に崩壊し，2000年代以降は
変動を伴いながら拡大して，近年は1,000兆円を超える規模に達している。

\subsubsection{東証の市場区分変更}

2022年4月4日，東京証券取引所の市場区分が変更された。従来の
「東証1部」「東証2部」「マザーズ」「ジャスダック」から，
「\textbf{プライム}」「\textbf{スタンダード}」「\textbf{グロース}」へ再編された。

変更の背景は次の2点である。
\begin{enumerate}
\item 市場区分のコンセプトの明確化。
\item 持続的な企業価値向上の動機付けとなる仕組み作り。
\end{enumerate}

新しい各市場区分のコンセプトは以下のとおりである。
\begin{itemize}
\item \textbf{プライム}：高い流動性と高いガバナンス水準(の企業向け)。
\item \textbf{スタンダード}：一定の流動性と基本的なガバナンス水準。
\item \textbf{グロース}：高い成長可能性をもつが，相対的にリスクは高い。
\end{itemize}

この市場区分変更に伴い，株価指数(TOPIX等)の算出方法も変更されている。

\subsection{流動性と証券化}

債券市場や株式市場は，証券に\textbf{流動性}を供給している。

\begin{definition}[流動性]
売買したいときに，適正な価格付近で，すぐに売買できるとき，\textbf{流動性が高い}
という。逆の状態を\textbf{流動性が低い}という。
\end{definition}

一方，銀行の個々の融資は市場で取引されない。すなわち流動性はない。仮に取引が
あっても外からは見えない(透明性を高める努力もあったが，十分ではない)。

\begin{definition}[証券化]
\textbf{証券化}とは，流動性のない資産を市場で売買できるようにする仕組みの一つ
である。多くの非流動性資産(ここでは融資)を集め，そこから生じるキャッシュ
フローをもとに，新たなリスク・リターン特性をもつ証券を組成し，販売する。
\end{definition}

2008年前後の金融危機で「活躍」したのは，住宅ローンの証券化商品である
RMBS (Residential Mortgage Backed Securities) である。さらに，
CDO (Collateralized Debt Obligations) や，合成CDOの素になる
CDS (Credit Default Swap) も金融危機において重要な役割を果たした。

\section*{参考：MBO(マネジメント・バイアウト)}

\textbf{マネジメント・バイアウト} (Management Buyout; MBO) とは，経営陣が自ら
株式を取得して株主になることをいう。自分達の経営戦略を円滑に進められるように
するための手法の一つであり，次のような場面で使われる。

\begin{itemize}
\item 経営陣による会社の買収(自己所有化)。
\item 企業グループ内の会社が，グループから分離独立するとき。
\item 他社による買収への対抗策。
\end{itemize}

上場企業が\textbf{敵対的買収}(買収者が，買収対象会社の取締役会の同意を得ないで
買収を仕掛けること)を回避するため，経営陣が MBO を行って市場から株式を
買い集め，上場を廃止して非公開会社にする，といったケースもある。
